

Relying on foundational LLMs alone typically reduces kernel development to a static, one-pass inference process. In contrast, LLM-based agents introduce autonomy and feedback into the optimization loop by enabling planning, tool use, and evaluation of intermediate results. This closed-loop, self-improving paradigm allows agent-based approaches to scale kernel optimization across diverse workloads and hardware platforms, while sustaining long-horizon, fatigue-free exploration. Concretely, we categorized recent agent-driven advancements into four structural dimensions: \emph{learning mechanisms}, \emph{external memory management}, \emph{hardware profiling integration}, and \emph{multi-agent orchestration}.

\subsection{Learning Mechanisms}

The first dimension of advancement concerns search strategies. Initial approaches view kernel generation as iterative refinement. Caesar (in \cite{ouyang2025kernelbench}) utilizes simple feedback loops to refine kernels, while Inference-Time Scaling~\cite{chen2025deepseekr1} demonstrates that scaling test-time compute and reflection significantly boost kernel quality. To manage complexity, PEAK~\cite{tariq2025peak} employs a stepwise, modular iterative refinement strategy, and ``Minimal Executable Programs'' is proposed~\cite{chu2025gpu} to enable efficient, isolated iteration without building costly full-scale applications. DiffAgent~\cite{zhu2026diffbench} adopts iterative refinement to accelerate diffusion models, TritonX~\cite{hammond2025agentic} uses iterative refinement within a state machine to cover kernels of complete PyTorch ATen backends, and KernelGen~\cite{kernelgen} leverages test-time scaling and reflection techniques to enable kernel generation for multi-chip backends. MaxCode~\cite{ou2026maxcode} further unifies existing iterative search methods under a max-reward reinforcement learning framework, combined with a natural language critique model converting raw execution feedback into diagnostic insights.


To escape local optima, recent frameworks adopt population-based evolution. Lange et al.~\cite{lange2025towards} optimize translation CUDA via mutation and crossover. FM Agent \cite{li2025fm} includes an evolutionary stage with the principles of diversity preservation, adaptive evolution, and multi-population dynamics. Advanced population dynamics are also introduced in EvoEngineer~\cite{Guo2025EvoEngineer}, which decouples traversal techniques from population management. GPU Kernel Scientist ~\cite{Andrews2025GPUKernelScientist} employs a multi-stage evolutionary workflow to address the challenge of optimizing HIP kernels for the AMD accelerators. And cuPilot~\cite{chen2025cupilot} guides evolution through high-level semantic strategies.


\subsection{External Memory Management}

Complex kernel optimization often requires domain-specific knowledge, such as CUDA APIs and hardware instruction sets that may be hallucinated or forgotten by the LLM. Agents in this category augment generation with external memory. The AI CUDA Engineer \cite{lange2025ai} leverages a vector database of high-quality kernel examples to ground the LLM’s generation, ensuring syntactic correctness and adherence to best practices in low-level programming. KernelEvolve~\cite{liao2025kernelevolvescalingagentickernel} further advances the external knowledge management paradigm by integrating a sophisticated hardware-specific knowledge base specifically tailored for heterogeneous AI accelerators.
Beyond retrieving unstructured textual context, recent work has explored utilizing structured representations as external memory to guide model inference. Work such as ReGraphT \cite{gong2025large} proposes a novel framework that treats a reasoning graph as a domain-specific external memory for CUDA code optimization. In this approach, the logical transitions between optimization states of large language models are externalized into a static, navigable graph structure for the small language model to retrieve.


\subsection{Hardware Profiling Integration}
% use what and how

The third dimension addresses the hardware-agnostic nature of standard LLMs by configuring the agent's persona profile with hardware specifications, and iteratively reasoning over performance profiling feedback. 

QiMent-TensorOp~\cite{zhang2025qimeng} triggers LLMs to analyze and distill low-level hardware documentation according to user input into the generation prompt, while QiMeng-GEMM~\cite{zhou2025qimenggemm} generates General Matrix Multiplication (GEMM) with the meta-prompt, which offers universal templates for various general optimization techniques and platform-specific optimization details. QiMeng-Attention~\cite{zhou2025qimeng} considers  target GPU architecture and instruction set to convert the high-level thinking language into low-level CUDA code, and  implements the high-performance FlashAttention on different GPUs. SwizzlePerf \cite{tschand2025swizzleperf} explicitly tackles the swizzling problem, which explicitly injects precise architectural specifications into the prompt context and restricts the search space specificaslly to swizzling patterns that focuses solely on maximizing the L2 cache hit rate.


Complementing this, agents leverage dynamic feedback. CUDA-LLM \cite{chen2025cuda} incorporates detailed target GPU specifications (e.g., warp size, cache size) into the agent's prompt. Simultaneously, compilation logs and runtime performance metrics are also aggregated to guide the optimization process. TritonForge~\cite{li2025tritonforge} utilizes profiling-guided feedback loops to iteratively analyze and identify performance bottlenecks. PRAGMA \cite{lei2025pragma} uses a specialized profiling module to parse low-level quantitative metrics into the interpretable natural language suggestion. KERNELBAND~\cite{ran2025kernelband} clusters runtime behavior of potential kernels to reduce the exploration space and utilizes profiling data  as context to guide the optimization strategies selection. 



\subsection{Multi-Agent Orchestration}


Recognizing that kernel development inherently involves heterogeneous skills ranging from algorithmic planning to low-level coding and debugging, recent works increasingly adopt multi-agent designs that explicitly decompose these responsibilities into coordinated roles.

STARK~\cite{dong2025stark} structures generation into Plan-Code-Debug phases to emulate human workflows, while AKG~\cite{du2025akgkernelagentmultiagent} leverages similar modularity to achieve cross-platform synthesis. Astra~\cite{wei2025astra} specializes this approach for production-grade SGLang kernels, focusing on tuning-focused agents. CudaForge~\cite{Zhang2025CudaForge} employs a Coder-Judge loop driven by hardware-level feedback, whereas KForge~\cite{sereda2025kforge} adapts this dual-agent model to new platforms using only single-shot example supervision. Addressing scale, KernelFalcon~\cite{kernelfalcon2024} employs a multi-agent system to tackle the challenge of GPU kernel generation of full machine learning architectures, where the system specifically addresses hierarchical task decomposition and delegation through coordinated manager and worker agents. Conversely, GEAK~\cite{Wang_2025_GEAK} targets AMD GPUs, integrating generation and reflection within a Triton-based workflow.

% STARK \cite{dong2025stark} exemplifies this paradigm by structuring kernel generation as a collaboration among Plan, Code, and Debug agents, closely mirroring a human engineering workflow. AKG \cite{du2025akgkernelagentmultiagent} similarly demonstrates the benefits of modularizing kernel generation into specialized sub-tasks handled by distinct agents, and achieves the cross-platform kernel synthesis. Other systems emphasize iterative generation–evaluation loops between complementary agent roles. CudaForge \cite{Zhang2025CudaForge} adopts a two-agent setup consisting of a coder agent and a judge agent, which repeatedly interact to generate and optimize CUDA kernels using hardware-level feedback, while KForge~\cite{sereda2025kforge} adopts an analogous dual-agent approach but significantly reduces the supervision requirement to a single-shot example when adapting to target new platforms. KernelFalcon~\cite{kernelfalcon2024} employs a multi-agent system to tackle the challenge of GPU kernel generation of full machine learning architectures, where the system specifically addresses hierarchical task decomposition and delegation through coordinated manager and worker agents.
% Astra \cite{wei2025astra} targets performance tuning of production kernels extracted from SGLang, orchestrating multiple specialized agents for planning, coding, testing, and profiling. In addition, GEAK \cite{Wang_2025_GEAK} developed for AMD GPUs, integrates agents for kernel generation, evaluation, reflection, and optimization within a Triton-based workflow, achieving substantial gains on rigorously designed benchmarks. 
